\section{関連研究}
\label{sec:related}

\emph{少数のターゲットラベルによるコーパス横断の感情認識．}
コーパス横断の評価は，音声感情認識で大きな性能の低下を以前から示しており，コーパスごとの正規化の動機となってきた\cite{schuller2010crosscorpus}．
少数のラベル付きターゲット発話を使う手法には，自発的な音声のためのシャムネットワークによる少数ショット学習\cite{feng2023fewshot}，ターゲットのラベルの量を増やしながら評価した敵対的なドメイン汎化\cite{gideon2021addog}，ソースコーパスでの事前の微調整（pre-finetuning）\cite{chen2023prefinetuning}，演技された音声・映像と即興の音声・映像の間の適応\cite{li2021acted}，凍結した音声特徴の上でロジスティック回帰を用いる能動学習\cite{xu2024crosstask}がある．
音声とテキストのモデルは，クラスごとに5から60本のラベルで\cite{he2025m4ser}，あるいは言語をまたいでコーパスごとに少数のラベルで\cite{zaidi2024enhancing}適応されている．
これらの研究は手法を提案するか評価するものであり，ターゲットのラベルが分類器のどの部分を変えるかを測った研究はなく，顔を含む構成はまだ調べられていない．
音声とテキストによるコーパス横断の認識の最近の研究も，コーパス横断のマルチモーダルの手法が少ないと指摘している\cite{ahn2025multitask}．
Gideonら\cite{gideon2021addog}は，いつ汎化し，いつ特化するかの定式化を今後の課題に挙げ，FengとChaspari\cite{feng2023fewshot}は，ソースの性質の影響とソースを組み合わせることの影響を挙げている．
本論文の判断基準と向きの分析と，ラベルの本数に換算したソースの価値は，線形ヘッドについてこれらの問いを扱う．
ラベルを使わない音声感情認識のテスト時適応では，エントロピー最小化と擬似ラベルがほとんどの設定で失敗すると報告されており\cite{dong2026test}，この結果は本論文の比較と一致する．

\emph{少数の標本から推定した線形の規則の誤差．}
標本線形判別の誤分類率の漸近展開は，次元によらない$1/N$の次数の項と，次元に比例する項を含み\cite{okamoto1963asymptotic}，Raudysは後者を実効次元で表した\cite{raudys1998intrinsic}．
本論文の判断基準のコスト$1/(2K)$と向きのコスト$q/K$は，罰則つきのロジスティック回帰についてこの分け方に従い，$q$には測った値を使う．
正解率の式\eqref{eq:twocloud}は，2つのガウス分布のクラスに対する線形の規則の標準的な式である\cite{desilva2023value,dobriban2018highdimensional}．
凍結した特徴と線形ヘッドについては，ドメイン調整回帰が超過リスクをターゲットの平均の誤差と残差で抑え\cite{rosenfeld2022domainadjusted}，分布のずれの下では微調整より線形プロービングが好ましいとされ\cite{kumar2022lpft}，正則化を0に近づける極限ではロジットへの加法的な補正が効果を失う\cite{kini2021labelimbalanced}．
少数のターゲット標本による転移の有限標本の解析は，回帰\cite{ha2026labeled}と1次元の判別\cite{desilva2023value}について存在する．
本論文は，実際のコーパスの対の上の多クラスロジスティック回帰について，対応する量を測る．

\emph{ターゲットの標本数で表したソースの価値．}
事前学習やソースの価値は，ターゲットのデータの実効的な量として文字数で\cite{hernandez2021scaling}，訓練例の数で\cite{lescao2021data}表され，また事前分布の実効標本サイズとして\cite{wiesenfarth2020quantification}，収束の速さの指数として\cite{hanneke2019value}表されてきた．
本論文の式\eqref{eq:value}は，推定した向きのコサインと標本サイズの古典的な関係\cite{raudys1998intrinsic,dobriban2018highdimensional}を逆に解いて，ソースの向きを感情ごとのラベルの本数で表し，本論文はそれをコーパスの対で測る．

\emph{合併と複数のソース．}
複数のドメインからの学習の上界では，ドメイン間の距離は罰則としてだけ働く\cite{bendavid2010theory,blitzer2007learning}．
合併した線形分類器の偏りは，高次元の極限で，タスクごとの切片をもつ中心化したマルチタスクのモデル\cite{tiomoko2021deciphering}と，切片をもたない合併の最小二乗法\cite{yang2025precise}について解かれており，ドメインの情報を明示的に加えると合併を上回りうる\cite{zhu2026domain}．
ニュートン法の1ステップは，点の集合の影響を近似する\cite{pregibon1981logistic,koh2019accuracy}．
これを，ターゲット専用の切片を設けずにターゲットの点を加える場合に当てはめると，式\eqref{eq:lambda}が得られる．
この式は，De Silvaら\cite{desilva2023value}の1次元の判断基準の移動を重みの経路で拡張したものであり，ターゲットがソースから分離できる度合いが，この経路を通じて判断基準を直す．

\emph{特権情報とモダリティ間の蒸留．}
タスクがラベルを共有し，入力のモダリティが異なるとき，追加のモダリティは特権情報（privileged information）である\cite{vapnik2009new,lopezpaz2016unifying}．
モダリティ間の蒸留は，対になったモダリティの間で教師信号を移す\cite{gupta2016cross}．
感情認識では，顔の先生がラベルのない動画で音声の生徒を学習させ\cite{albanie2018emotion}，音声とテキストを融合した先生が音声の生徒を学習させ\cite{srinivasan2022representation}，音声と映像の先生や映像の先生が，ソフトラベルや埋め込みを音声の生徒に渡す\cite{pendyala2025leveraging,hajavi2024audio}．
単一モダリティの先生は融合の生徒を改善しうるし\cite{du2021improving,xue2021multimodal,rakib2025g2d}，融合の出力で単一モダリティのヘッドを学習させることもできる\cite{jaritz2020xmuda}．
これらの研究では，先生をターゲットのコーパスか第3のコーパスで学習するか，生徒が推論時にすべてのモダリティを使う．
本論文の先生はソースコーパスだけで学習した線形分類器であり，学び直さずにターゲットのプールにラベルを付ける．
生徒は，ソースから作った出発点（start）にアンカーした単一モダリティの線形ヘッドである．

\emph{弱いラベル，クラスの釣り合い，アンカー，選択．}
固定した先生が，クラスの異なるタスクをまたぐ表現学習のために\cite{phoo2021selftraining}，また固定した重みによるトランスダクティブな少数ショット分類のために\cite{zanella2024boosting}，ラベルのないデータすべてにソフトラベルを付けた研究があり，モダリティごとの先生をドメインをまたいで単一モダリティに蒸留した研究もある\cite{hatano2024multimodal}．
擬似ラベルをクラス間で釣り合わせることは一般的であるが\cite{huang2022unsupervised,xie2020selftraining,hu2021leveraging,berthelot2022adamatch}，ラベルのない集合のクラスが偏っていると，その釣り合わせが害になりうる\cite{veilleux2021realistic}．
2次の罰則でゼロショット分類器にアンカーした線形プローブ\cite{silvarodriguez2024closer}と仮説転移（hypothesis transfer）\cite{kuzborskij2017fast,li2018explicit}は，重みを事前の解にアンカーし，1つのドメインを除く検証（leave-one-domain-out）はターゲットのラベルなしでハイパーパラメータを選ぶ\cite{gulrajani2021domainbed}．
提案法はこれらの構成要素を組み合わせる．
提案法に固有なのは，ソースだけで学習した別のモダリティの先生，ソースとターゲットが感情の集合を共有するので分類器そのものの学習に使う弱いラベル，重みと切片に別々に置くアンカー，先生と弱いラベルの重みを$K$ごとに選ぶことである．

\emph{弱いラベルによる利得とラベルを使わない選択の理論．}
弱い先生のラベルで学習した生徒の利得は，両者の不一致（misfit）で下から抑えられ\cite{charikar2024quantifying}，両者の表現の食い違いと関係づけられている\cite{dong2025discrepancies,xue2025representations}．
自己訓練は，徐々に進む分布のずれの下で解析されている\cite{kumar2020understanding}．
ガウスの線形モデルでは，最適な計画は観測された応答によらず\cite{chaloner1984optimal,chaloner1995bayesian}，予算が小さいときはランダムな選択や典型的な点の選択を上回るのは難しい\cite{hacohen2022active}．
命題\ref{prop:span}は表現定理\cite{kimeldorf1971results,scholkopf2001generalized}から導かれる．
本論文はこれらの結果を，出発点にアンカーした生徒に用いる．
先生の出発点に対するUARの優位が利得の規模を決め，ラベルを使わない選択は，修正される部分空間を出発点の向きの誤差にそろえられない．
